\subsection{Results}

For running the application the following commands were used for compiling,
uploading the project and monitoring the port to observe the periodic reports.

\begin{verbatim}
# to find to which port number the board is connected
arduino-cli board list
# compilation of the program by indicating the target hardware
arduino-cli compile --fqbn esp32:esp32:esp32 sketch
# uploading the program executables to the board
arduino-cli upload -p /dev/cu.usbserial-portnumber --fqbn esp32:esp32:esp32 sketch
# monitoring the port at 115200 baud (ESP32 default)
arduino-cli monitor -p /dev/cu.usbserial-portnumber --config baudrate=115200
\end{verbatim}

Below is the output of the commands listed above

\begin{verbatim}
$ arduino-cli compile --fqbn esp32:esp32:esp32 sketch
Sketch uses 267498 bytes (20%) of program storage space.
Global variables use 21516 bytes (6%) of dynamic memory,
leaving 306164 bytes for local variables.

$ arduino-cli upload -p /dev/cu.usbserial-0001 --fqbn esp32:esp32:esp32 sketch
New upload port: /dev/cu.usbserial-0001 (serial)

$ arduino-cli monitor -p /dev/cu.usbserial-0001 --config baudrate=115200
Using generic monitor configuration.

Monitor port settings:
  baudrate=115200

Connecting to /dev/cu.usbserial-0001. Press CTRL-C to exit.

Lab 4 - FreeRTOS Button Press Monitor
Reports every 10 seconds
--- Report ---
Total presses: 5
Short (<500ms): 3
Long  (>=500ms): 2
Avg duration: 487 ms
--------------
--- Report ---
Total presses: 0
Short (<500ms): 0
Long  (>=500ms): 0
Avg duration: N/A
--------------
\end{verbatim}

The system behavior is as follows:
\begin{itemize}
  \item When a short press ($< 500$ ms) is detected, the \textbf{green LED} turns on and the
    red LED turns off. Task 1 gives the binary semaphore and Task 2 picks it up on its next
    50\,ms tick. The \textbf{yellow LED} blinks 5 times rapidly.
  \item When a long press ($\geq 500$ ms) is detected, the \textbf{red LED} turns on and the
    green LED turns off. The \textbf{yellow LED} blinks 10 times rapidly.
  \item Every 10 seconds, Task 3 atomically snapshots and resets statistics under the mutex,
    then prints the report to the serial monitor showing the total number of presses, the count
    of short and long presses, and the average press duration.
  \item When no presses occur during a reporting interval, the report shows zeros and
    ``N/A'' for the average duration.
  \item The monitor must be opened at \textbf{115200 baud} to receive readable output, as the
    ESP32 serial interface operates at this rate.
\end{itemize}

Below there is the visual representation of the circuit

\begin{figure}[h]
    \centering
    \includegraphics[width=0.6\textwidth]{img/ledoff.jpg}
    \caption{Circuit in idle state}
    \label{fig:ledoff}
\end{figure}

Figure \ref{fig:ledoff} shows the circuit in its idle state with all LEDs off.

A video demonstrating how the program works is posted on youtube hosting platform.
Video can be accessed using this link \href{https://youtu.be/teMOPxAjXb4}{https://youtu.be/teMOPxAjXb4}
